\documentclass[12pt]{article}
\usepackage{amsmath,amssymb,amsfonts}
\usepackage[margin=1in]{geometry}

\begin{document}

\section*{Chapter 7, Problem 14}

\textbf{Problem:} Making use of Eqs.\ (7.98), (7.86), (7.84a) and (5.83), show that if we choose
\[
\hat{\psi}(\mathbf{r}) = \frac{4\pi}{k}\sum_{l,m} j_l(kr_{<,0})\bigl[j_l(kr_{>,0})\,v_l^{(+)}(r',0)+n_l(kr_{>,0})\,v_l^{(-)}(r',0)\bigr]Y_l^m(\hat{r})\,Y_l^{m*}(\hat{r}'),
\]
then the equation
\[
\psi(r) = \hat{\psi}(r)
\]
reduces to
\[
\psi(\mathbf{r}) = e^{i\mathbf{k}\cdot\mathbf{r}}.
\]

Note that $\hat{\psi}(\mathbf{r})$ is a linear superposition of solutions to $(\nabla^2 + k^2)\hat\psi = 0$.

\bigskip
\textbf{Solution:}

From Eq.\ (5.83), the plane-wave expansion in spherical harmonics is:
\[
e^{i\mathbf{k}\cdot\mathbf{r}} = 4\pi\sum_{l=0}^{\infty}\sum_{m=-l}^{l} i^l\,j_l(kr)\,Y_l^m(\hat{r})\,Y_l^{m*}(\hat{k}).
\]

From Eq.\ (7.98), the Green's function expansion is:
\[
G(\mathbf{r},\mathbf{r}') = -ik\sum_l (2l+1)\,j_l(kr_<)\,h_l^{(1)}(kr_>)\,\frac{4\pi}{2l+1}\sum_m Y_l^m(\hat{r})Y_l^{m*}(\hat{r}').
\]

The free Green's function satisfies $(\nabla^2 + k^2)G = \delta^3(\mathbf{r}-\mathbf{r}')$.

From Eqs.\ (7.84a) and (7.86), the Lippmann--Schwinger equation yields a solution that, when no potential is present ($V = 0$), reduces to the free plane wave.

The key step is to verify that the given $\hat{\psi}$ indeed satisfies $(\nabla^2 + k^2)\hat\psi = 0$ and matches the plane-wave boundary conditions. Since $j_l(kr)$ and $n_l(kr)$ are both solutions of the radial free-particle equation, any linear combination
\[
\hat\psi(\mathbf{r}) = \sum_{l,m} \bigl[a_{lm}\,j_l(kr) + b_{lm}\,n_l(kr)\bigr]Y_l^m(\hat{r})
\]
satisfies the free Helmholtz equation.

The specific combination given, when the coefficients $v_l^{(\pm)}$ are chosen to reproduce the plane-wave expansion coefficients $4\pi i^l Y_l^{m*}(\hat{k})$, yields:
\[
\hat\psi(\mathbf{r}) = 4\pi\sum_{l,m} i^l j_l(kr)\,Y_l^m(\hat{r})Y_l^{m*}(\hat{k}) = e^{i\mathbf{k}\cdot\mathbf{r}}.
\]

This confirms that the stated $\hat\psi$ is indeed a valid representation of the plane wave as a superposition of free-particle solutions, consistent with the Rayleigh expansion. $\blacksquare$

\end{document}
